\documentclass[11pt,reqno]{amsart}

\usepackage[T1]{fontenc}
\usepackage{newtxtext,newtxmath}
\let\Bbbk\relax
\usepackage{amsmath,amssymb,amsfonts}
\usepackage{microtype}
\usepackage{xcolor}
\usepackage[numbers]{natbib}
\usepackage[colorlinks=true,linkcolor=blue,citecolor=blue!60!black,urlcolor=blue!60!black]{hyperref}

\begin{document}

\title{Dynamics of representation, learning, and routing in deep networks}
\maketitle

\section{Literature}

Deep neural networks are implemented as compositions of finite-dimensional maps whose parameters are adjusted by gradient-based optimization. Residual networks made the relation between depth and discrete-time state evolution particularly explicit \cite{he2016resnet}, while neural ordinary differential equations formalized the corresponding continuous-depth limit \cite{chen2018neuralode}. The training dynamics have been studied from several complementary limits. Exact analyses of deep linear networks isolate mode-dependent learning times \cite{saxe2014}, the neural tangent kernel describes an infinite-width regime in which the feature map is effectively frozen \cite{jacot2018}, and mean-field limits instead retain a distribution of trainable units whose evolution can be written as a transport or Wasserstein gradient flow \cite{chizat2018}.

Transformers add state-dependent interactions between the elements of a sequence. The original architecture defines attention through learned query, key, and value projections followed by a softmax normalization \cite{vaswani2017}. Several later theories recover familiar dynamical objects from this computation. Modern Hopfield networks identify a particular attention update with an associative-memory fixed-point map generated by an energy \cite{ramsauer2021}. Mean-field descriptions treat tokens as interacting particles and derive continuum transport equations for the empirical token distribution \cite{castin2025,rigollet2025}. A recent formulation also couples the token distribution evolving through depth to a parameter distribution evolving through training, placing forward propagation and learning in one mean-field system \cite{herty2026}.

Mixture-of-experts architectures add input-dependent allocation of computation. Sparse gating was introduced as a way to increase model capacity while evaluating only a small subset of experts \cite{shazeer2017}. Switch Transformers use top-one routing together with auxiliary balancing mechanisms \cite{fedus2022}, while expert-choice routing reverses the assignment direction and lets experts select tokens subject to capacity constraints \cite{zhou2022}. Current theory has begun to analyze joint router--expert feature learning \cite{liao2026}, while empirical work shows that expert specialization can be fine grained and superposed rather than aligned with broad semantic domains \cite{lasy2026}. These results motivate a common notation that keeps the implemented computation visible while distinguishing exact identities, limiting descriptions, and physical interpretations.

The purpose here is organizational. We write the forward pass, optimizer, attention mechanism, and expert routing in a single notation and then state the dynamical descriptions that follow from them. The resulting objects are the representation state, the trainable parameters, and the allocation weights. No general energy principle or evolutionary interpretation is assumed unless it follows from the specific model under consideration.

\section{Results}

\subsection{From text to a Transformer block}

Consider a sequence of text converted by a tokenizer into integer token identifiers
\begin{equation}
(u_1,\ldots,u_N),
\qquad
u_i\in\{1,\ldots,V\},
\end{equation}
where $V$ is the vocabulary size. Tokenization is a discrete preprocessing map; the neural network starts from the corresponding embedding vectors. Let $\mathbf E\in\mathbb R^{V\times d}$ be the embedding matrix and let $\mathbf p_i\in\mathbb R^d$ denote positional information. The initial representation of token $i$ is
\begin{equation}
\mathbf x_i^{0}
=
\mathbf E_{u_i}
+
\mathbf p_i.
\label{eq:embedding}
\end{equation}
Thus a sequence is represented inside the network by $N$ vectors $\mathbf x_i^\ell\in\mathbb R^d$, with the superscript $\ell$ denoting the layer.

For one attention head, queries, keys, and values are linear projections of the current representations,
\begin{equation}
\mathbf q_i
=
\mathbf W_Q\mathbf x_i,
\qquad
\mathbf k_i
=
\mathbf W_K\mathbf x_i,
\qquad
\mathbf v_i
=
\mathbf W_V\mathbf x_i.
\label{eq:qkv}
\end{equation}
The query $\mathbf q_i$ specifies what token $i$ uses to score other tokens, the key $\mathbf k_j$ specifies how token $j$ is scored, and the value $\mathbf v_j$ is the vector transmitted if that interaction receives weight. With key dimension $d_k$, the score from token $i$ to token $j$ is
\begin{equation}
s_{ij}
=
\frac{
\mathbf q_i^{\mathsf T}\mathbf k_j
}{
\sqrt{d_k}
}
+
m_{ij},
\label{eq:score}
\end{equation}
where $m_{ij}$ is an optional mask. In a causal decoder,
\begin{equation}
m_{ij}
=
\begin{cases}
0, & j\leq i,\\
-\infty, & j>i,
\end{cases}
\end{equation}
so a token cannot attend to future positions.

The softmax maps the score vector of token $i$ to the probability simplex,
\begin{equation}
a_{ij}
=
\frac{
\exp(s_{ij})
}{
\sum_{r=1}^{N}\exp(s_{ir})
},
\qquad
a_{ij}\geq0,
\qquad
\sum_j a_{ij}=1.
\label{eq:attention-softmax}
\end{equation}
The output of the head is the weighted average
\begin{equation}
\mathbf y_i
=
\sum_{j=1}^{N}
a_{ij}\mathbf v_j.
\label{eq:attention-output}
\end{equation}
Equations~\eqref{eq:qkv}--\eqref{eq:attention-output} are the basic self-attention computation. The interaction coefficients are not fixed couplings: they are recomputed from the current representations at every layer.

With $H$ heads, each head has its own projections and produces $\mathbf y_i^{(h)}$. Their outputs are concatenated and projected,
\begin{equation}
\operatorname{MHA}(\mathbf x)_i
=
\mathbf W_O
\begin{pmatrix}
\mathbf y_i^{(1)}\\
\vdots\\
\mathbf y_i^{(H)}
\end{pmatrix}.
\end{equation}
A pre-normalization Transformer block can then be written as
\begin{align}
\mathbf x_i'
&=
\mathbf x_i^\ell
+
\operatorname{MHA}
\left(
\operatorname{Norm}(\mathbf x^\ell)
\right)_i,
\label{eq:attention-residual}
\\
\mathbf x_i^{\ell+1}
&=
\mathbf x_i'
+
\operatorname{MLP}
\left(
\operatorname{Norm}(\mathbf x_i')
\right),
\label{eq:mlp-residual}
\end{align}
where the feed-forward sublayer is typically
\begin{equation}
\operatorname{MLP}(\mathbf x)
=
\mathbf W_2
\sigma
\left(
\mathbf W_1\mathbf x+\mathbf b_1
\right)
+
\mathbf b_2.
\end{equation}
The exact activation, normalization, positional encoding, and residual arrangement depend on the architecture. Equations~\eqref{eq:attention-residual} and \eqref{eq:mlp-residual} isolate the common residual structure needed below.

For autoregressive language modeling, the final representation at position $i$ is projected to vocabulary logits,
\begin{equation}
r_{ik}
=
\left(
\mathbf W_{\mathrm{out}}\mathbf x_i^L
+
\mathbf b_{\mathrm{out}}
\right)_k.
\end{equation}
The model probability for the next token is another softmax,
\begin{equation}
p_{ik}
=
\frac{
\exp(r_{ik})
}{
\sum_{q=1}^{V}\exp(r_{iq})
}.
\label{eq:output-softmax}
\end{equation}
For target token $y_i$, the local cross-entropy loss is $-\log p_{i y_i}$. The same mathematical map therefore appears both inside attention and at the output, but over different index sets and with different meanings.

\subsection{Parameter learning and its continuous approximation}

Collect all trainable parameters into $\theta$. For a minibatch $B_n$, backpropagation computes
\begin{equation}
\mathbf g_n
=
\nabla_\theta L_{B_n}(\theta_n).
\end{equation}
The simplest update is stochastic gradient descent,
\begin{equation}
\theta_{n+1}
=
\theta_n
-
\eta\mathbf g_n.
\label{eq:sgd}
\end{equation}
This discrete update is the implemented object from which the usual gradient-flow description is obtained. If the step size is small and the stochasticity of minibatching is neglected at the scale of interest, Eq.~\eqref{eq:sgd} gives
\begin{equation}
\dot\theta
=
-
\nabla_\theta L(\theta).
\label{eq:gradient-flow}
\end{equation}
Equation~\eqref{eq:gradient-flow} is therefore a continuous approximation, not a replacement for the optimizer used in training.

Large Transformers are commonly trained with adaptive optimizers such as AdamW. Ignoring bias-correction notation for compactness, its additional state can be written as
\begin{align}
\mathbf m_n
&=
\beta_1\mathbf m_{n-1}
+
(1-\beta_1)\mathbf g_n,
\\
\mathbf v_n
&=
\beta_2\mathbf v_{n-1}
+
(1-\beta_2)\mathbf g_n^2,
\\
\theta_{n+1}
&=
\theta_n
-
\eta
\frac{
\widehat{\mathbf m}_n
}{
\sqrt{\widehat{\mathbf v}_n}+\epsilon
}
-
\eta\lambda\theta_n.
\label{eq:adamw}
\end{align}
The optimizer is then itself a dynamical system with memory variables $\mathbf m_n$ and $\mathbf v_n$. A Euclidean gradient flow such as Eq.~\eqref{eq:gradient-flow} should only be used when the corresponding reduction is justified.

\subsection{Depth as state dynamics}

The residual structure gives a direct route from the implemented network to a continuous state equation. Introduce an explicit layer step $\Delta\tau$ and write
\begin{equation}
\mathbf x_i^{\ell+1}
=
\mathbf x_i^\ell
+
\Delta\tau
F_i
\left(
\mathbf x^\ell;\theta
\right).
\label{eq:residual-step}
\end{equation}
If the layers sample a sufficiently regular vector field as $\Delta\tau\to0$, the continuous-depth limit is
\begin{equation}
\partial_\tau\mathbf x_i
=
F_i
\left(
\mathbf x;\theta
\right).
\label{eq:depth-flow}
\end{equation}
The variable $\tau$ describes depth or inference time. It is distinct from the training time that parametrizes Eq.~\eqref{eq:gradient-flow}.

Attention makes Eq.~\eqref{eq:depth-flow} an interacting system because $F_i$ depends on the complete token set through the coefficients $a_{ij}$. At finite $N$, the empirical token measure is
\begin{equation}
\mu_\tau^N
=
\frac{1}{N}
\sum_{i=1}^{N}
\delta_{\mathbf x_i(\tau)}.
\label{eq:empirical-token}
\end{equation}
Under mean-field assumptions, the limit $N\to\infty$ can be described by a transport equation of the form
\begin{equation}
\partial_\tau\mu
+
\nabla_{\mathbf x}\cdot
\left(
\mu\,
\mathbf v[\mu]
\right)
=
0.
\label{eq:token-pde}
\end{equation}
For specific attention variants, this limit and its well-posedness can be derived explicitly \cite{castin2025,rigollet2025}. The finite Transformer and the continuum equation are therefore different mathematical levels: the particle system is the architecture-level object, while Eq.~\eqref{eq:token-pde} is a large-token reduction.

A similar distinction appears in parameter space. In suitable infinite-width or infinite-head limits, a collection of trainable units can be represented by a parameter measure $\rho_t$. Mean-field learning then takes the generic Wasserstein form
\begin{equation}
\partial_t\rho
=
\nabla_\theta\cdot
\left[
\rho
\nabla_\theta
\frac{\delta L}{\delta\rho}
\right].
\label{eq:parameter-pde}
\end{equation}
Recent Transformer theory couples a token distribution evolving through depth to a parameter distribution evolving through training \cite{herty2026}. This makes the separation between representation dynamics and learning dynamics explicit even when both are described by continuum equations.

\subsection{Energy and fixed-point descriptions}

An energy description is available for particular attention-like systems, but not for an arbitrary Transformer block. Modern Hopfield networks provide a useful exact case. Let $\mathbf X=(\boldsymbol\xi_1,\ldots,\boldsymbol\xi_M)$ contain stored patterns and define
\begin{equation}
E(\mathbf x)
=
\frac{1}{2}
\mathbf x^{\mathsf T}\mathbf x
-
\frac{1}{\beta}
\log
\sum_{\mu=1}^{M}
\exp
\left(
\beta
\boldsymbol\xi_\mu^{\mathsf T}\mathbf x
\right).
\label{eq:hopfield-energy}
\end{equation}
Terms independent of $\mathbf x$ have been omitted. Its gradient is
\begin{equation}
\nabla E(\mathbf x)
=
\mathbf x
-
\mathbf X
\operatorname{softmax}
\left(
\beta\mathbf X^{\mathsf T}\mathbf x
\right).
\end{equation}
The associated fixed-point update,
\begin{equation}
\mathbf x^{+}
=
\mathbf X
\operatorname{softmax}
\left(
\beta\mathbf X^{\mathsf T}\mathbf x
\right),
\label{eq:hopfield-update}
\end{equation}
has the same query--key--value algebra as an attention head under the corresponding identification of patterns and projections \cite{ramsauer2021}. Fixed points satisfy $\mathbf x^+=\mathbf x$. This connection supports an energetic interpretation for the specified construction; it does not imply that the complete trained Transformer admits one scalar energy whose gradient generates every layer.

\subsection{Mixture-of-experts as allocation of computation}

A mixture-of-experts layer replaces one feed-forward map by $E$ expert maps $F_e(\mathbf x;\theta_e)$ together with a router. For a linear router,
\begin{equation}
z_{ie}
=
\mathbf w_e^{\mathsf T}\mathbf x_i
+
b_e
\end{equation}
is the logit assigning token $i$ to expert $e$. Dense routing uses
\begin{equation}
a_{ie}
=
\frac{
\exp(z_{ie})
}{
\sum_{f=1}^{E}\exp(z_{if})
},
\qquad
\sum_e a_{ie}=1,
\label{eq:router-softmax}
\end{equation}
and returns
\begin{equation}
F_{\mathrm{MoE}}(\mathbf x_i)
=
\sum_{e=1}^{E}
a_{ie}
F_e(\mathbf x_i;\theta_e).
\label{eq:dense-moe}
\end{equation}

Sparse routing first chooses a support
\begin{equation}
S_i
=
\operatorname{TopK}
\left(
z_{i1},\ldots,z_{iE}
\right)
\end{equation}
and then normalizes only over the selected experts,
\begin{equation}
\widetilde a_{ie}
=
\begin{cases}
\dfrac{
\exp(z_{ie})
}{
\sum_{f\in S_i}\exp(z_{if})
},
&
e\in S_i,
\\[8pt]
0,
&
e\notin S_i.
\end{cases}
\label{eq:topk}
\end{equation}
The output is obtained from Eq.~\eqref{eq:dense-moe} after replacing $a_{ie}$ by $\widetilde a_{ie}$. The top-$k$ operation is important mathematically because the routing map is smooth while $S_i$ is fixed but changes non-smoothly when two logits exchange order.

\subsection{Softmax induces a replicator equation}

The relation between softmax routing and replicator dynamics follows directly from differentiation. Consider one token and omit its index. Let
\begin{equation}
a_e
=
\frac{
\exp(z_e)
}{
\sum_f\exp(z_f)
}
\end{equation}
and let the logits follow any differentiable path $z_e(t)$. Define
\begin{equation}
f_e
=
\dot z_e,
\qquad
\overline f
=
\sum_f a_f f_f.
\end{equation}
Differentiating the softmax gives the exact identity
\begin{equation}
\dot a_e
=
a_e
\left(
f_e-\overline f
\right).
\label{eq:replicator}
\end{equation}
Thus smooth softmax weights evolve in replicator form when viewed along the evolution of their logits. No evolutionary assumption is needed for Eq.~\eqref{eq:replicator}; it is a coordinate identity on the simplex.

For a router $z_e=z_e(\mathbf x;\phi)$ with router parameters $\phi$, the induced rate is
\begin{equation}
f_e
=
\nabla_{\mathbf x}z_e\cdot\dot{\mathbf x}
+
\nabla_\phi z_e\cdot\dot\phi.
\label{eq:router-rate}
\end{equation}
If the representation is held fixed during the parameter update, only the second term remains. If $\phi$ follows gradient descent, then $\dot\phi=-\nabla_\phi L$ and Eq.~\eqref{eq:router-rate} connects optimizer dynamics to an induced simplex dynamics. Interpreting $f_e$ as fitness is an additional modeling step. For top-$k$ routing, Eq.~\eqref{eq:replicator} holds piecewise on a fixed active support, while support changes introduce switching events.

\subsection{Specialization and collapse are distinct observables}

Let $X$ denote a token drawn from the data distribution and let $a_e(X)$ be its soft routing weights. Define the average occupation of expert $e$ by
\begin{equation}
q_e
=
\mathbb E_X
\left[
a_e(X)
\right].
\end{equation}
The global load entropy is
\begin{equation}
H(q)
=
-
\sum_e q_e\log q_e,
\end{equation}
while the mean routing entropy is
\begin{equation}
U
=
\mathbb E_X
\left[
-
\sum_e
a_e(X)\log a_e(X)
\right].
\end{equation}
Their difference is
\begin{equation}
S
=
H(q)-U
=
I(X;E),
\label{eq:specialization}
\end{equation}
where $E$ is the random expert selected according to $a(X)$. Equation~\eqref{eq:specialization} is the mutual information between tokens and expert assignments. It measures token-dependent specialization independently of whether the global load is balanced. A complementary collapse measure is
\begin{equation}
C
=
\log E-H(q).
\label{eq:collapse}
\end{equation}
Uniform soft routing has $C=0$ and $S=0$. Balanced deterministic routing has $C=0$ and large $S$. Routing concentrated on one expert has large $C$ and $S=0$. These cases separate load imbalance from specialization without assuming that an expert corresponds to one semantic domain.

\subsection{A common dynamical notation}

The implemented system can now be written with two distinct time variables. Let $\tau$ denote depth and $t$ denote training time. The representation state is $\mathbf x_i(\tau,t)$, the ordinary network parameters are $\theta(t)$, and router parameters are $\phi(t)$. A minimal continuous notation is
\begin{align}
\partial_\tau\mathbf x_i
&=
F_i
\left(
\mathbf x;\theta,\phi
\right),
\label{eq:common-state}
\\
\dot\theta
&=
G_\theta
\left(
\theta,\phi
\right),
\label{eq:common-weights}
\\
\dot\phi
&=
G_\phi
\left(
\theta,\phi
\right),
\label{eq:common-router}
\\
a_{ie}
&=
\operatorname{softmax}
\left[
z_i(\mathbf x_i;\phi)
\right]_e.
\label{eq:common-allocation}
\end{align}
For ordinary gradient flow, $G_\theta=-\nabla_\theta L$ and $G_\phi=-\nabla_\phi L$. For another optimizer, its internal state must be added to Eqs.~\eqref{eq:common-weights} and \eqref{eq:common-router}. The allocation $a_{ie}$ is normally an instantaneous function of the representation and router parameters, rather than an independent state variable. Its dynamics follow from Eq.~\eqref{eq:replicator} when the softmax is differentiable.

The main theoretical descriptions fit into this notation at different levels:
\begin{center}
\begin{tabular}{lll}
\hline
relation & status & object \\
\hline
residual layer $\to$ ODE & continuous-depth limit & $\mathbf x$ \\
finite attention $\to$ particles & exact rewriting & $\{\mathbf x_i\}$ \\
particles $\to$ transport PDE & mean-field limit & $\mu$ \\
SGD $\to$ gradient flow & small-step approximation & $\theta$ \\
wide model $\to$ parameter PDE & mean-field limit & $\rho$ \\
Hopfield update $\leftrightarrow$ attention & exact for matched form & $\mathbf x$ \\
softmax path $\to$ replicator equation & exact identity & $a$ \\
specialization $\leftrightarrow$ niche formation & interpretation & $a,\theta$ \\
\hline
\end{tabular}
\end{center}

This separation prevents formally different statements from being merged under the same physical analogy. The finite network, its continuum limits, and the interpretation of the resulting variables remain identifiable at every step.

\subsection{Fixed points and coupled questions}

Each level has its own stationary condition. A representation fixed point satisfies
\begin{equation}
F_i(\mathbf x^*;\theta,\phi)=0,
\end{equation}
while a stationary point of Euclidean gradient flow satisfies
\begin{equation}
\nabla_\theta L(\theta^*,\phi^*)=0,
\qquad
\nabla_\phi L(\theta^*,\phi^*)=0.
\end{equation}
For the induced allocation dynamics in Eq.~\eqref{eq:replicator},
\begin{equation}
a_e^*>0
\quad\Longrightarrow\quad
f_e^*=\overline f^*
\end{equation}
at an interior stationary allocation. Boundary equilibria allow inactive experts with $a_e^*=0$.

The coupled problem is therefore not one fixed-point problem but the interaction of several dynamical levels. Relevant control parameters include depth, learning rate, router temperature, expert capacity, sparsity $k$, width, and explicit load-balancing terms. Depending on the approximation, one can then study stability, symmetry breaking between initially equivalent experts, switching of top-$k$ supports, and timescale separation between representation and parameter evolution. These are questions posed by the common formulation rather than additional assumptions about how every Transformer must behave.

\section{Discussion}

The same implementation can support several mathematical descriptions, but their status differs. Queries, keys, values, softmax attention, residual updates, optimizer states, and expert routing are finite computations. Continuous depth, mean-field token equations, parameter transport equations, and diffusion approximations require limiting assumptions. Energy functions exist for particular attention-like constructions, while evolutionary language becomes literal only after specifying a dynamics for which the simplex rates can be interpreted as fitness differences.

The common notation keeps these layers connected without identifying them. The representation $\mathbf x$ describes what the network computes, the parameters $\theta$ and $\phi$ describe what training changes, and the routing weights $a$ describe how computation is allocated. Existing theories can then be read as reductions of this larger discrete system rather than as competing descriptions of unrelated objects.

\bibliographystyle{unsrtnat}
\bibliography{refs}

\end{document}
